\subsection{Relation to the null-like warped \texorpdfstring{$SL(2)$}{SL(2)} Landau--Lifshitz pair}
\label{sec:class5-known-null-pair}

The Class 5 pair can be compared directly with the null-like warped $SL(2)$
Landau--Lifshitz pair of Ref.~\cite{KameyamaYoshida2014}.  Let
$n^0,n^1,n^2$ denote the fields of that model and define
\begin{equation}
(n^0)^2-(n^1)^2-(n^2)^2=1,
\qquad
n^{\pm}:=\frac{n^0\pm n^1}{\sqrt2}.
\label{eq:class5-known-n-fields}
\end{equation}
For the auxiliary-space basis, set
\begin{equation}
T^0=\frac{\ii}{2}\sigma^2,
\qquad
T^1=\frac12\sigma^1,
\qquad
T^2=\frac12\sigma^3,
\qquad
T^{\pm}=\frac{T^0\pm T^1}{\sqrt2}.
\label{eq:class5-known-generators}
\end{equation}
Introduce a nonzero normalization $A_K$, let $B_K:=-2/A_K$, and denote the
null-like deformation coupling and spectral parameter by $C_K$ and $z_K$.
The Lax pair of Ref.~\cite{KameyamaYoshida2014}, with these constants renamed,
is
\begin{align}
U_K
={}&\frac{A_K}{z_K}
\left[
-\left(n^+-\frac{C_Kz_K^2}{2}n^-\right)T^-
+n^2T^2-n^-T^+
\right],
\label{eq:class5-known-UK}\\
V_K
={}&\frac{B_K}{z_K}\Biggl[
-\left(
 n^2\partial_xn^+-n^+\partial_xn^2
 -\frac{C_Kz_K^2}{2}
 (n^-\partial_xn^2-n^2\partial_xn^-)
\right)T^-
\nonumber\\
&\hspace{17mm}
+(n^-\partial_xn^+-n^+\partial_xn^-)T^2
-(n^-\partial_xn^2-n^2\partial_xn^-)T^+
\Biggr]
\nonumber\\
&+\frac{A_KB_K}{z_K^2}
\left[
-\left(n^++\frac{C_Kz_K^2}{2}n^-\right)T^-
+n^2T^2-n^-T^+
\right].
\label{eq:class5-known-VK}
\end{align}

Define a rotated spin $\widetilde{\boldsymbol S}$ by
\begin{equation}
\boldsymbol S(x,t)
=\mathcal R(x)\widetilde{\boldsymbol S}(x,t),
\qquad
\mathcal R(x):=
\exp\!\left(-\frac{\alpha_5x}{2}M_5\right).
\label{eq:class5-known-field-rotation}
\end{equation}
Since $M_5^{\mathsf T}=-M_5$, this transformation preserves
$\boldsymbol S^{\mathsf T}\boldsymbol S=1$.  The fields and parameters are
related by
\begin{equation}
\begin{gathered}
\begin{pmatrix}n^0(x,t_K)\\ n^1(x,t_K)\\ n^2(x,t_K)\end{pmatrix}
=
\operatorname{diag}(1,-\ii,-\ii)\,
\widetilde{\boldsymbol S}(x,t_K/A_K^2),
\\[1mm]
C_K=\frac{\alpha_5^2}{2A_K^2},
\qquad
z_K=2\ii A_K\lambda,
\qquad
t_K=A_K^2t.
\end{gathered}
\label{eq:class5-known-dictionary}
\end{equation}
Equation~\eqref{eq:class5-known-dictionary} maps the complexified spin spaces;
the two real sections are different.

For the auxiliary-space transformation, let
\begin{equation}
Q=\begin{pmatrix}0&\ii\\1&0\end{pmatrix},
\qquad
G(x,\lambda)
=\Id_2+\alpha_5\left(\lambda-\frac{x}{2}\right)E_{12},
\qquad
\mathcal G:=GQ^{-1},
\label{eq:class5-known-gauge}
\end{equation}
where $E_{12}$ is the matrix unit with only its $(1,2)$ entry nonzero.
Define
\begin{equation}
U_{5,0}:=U_5-\frac12\tr U_5\,\Id_2,
\qquad
V_{5,0}:=V_5-\frac12\tr V_5\,\Id_2.
\label{eq:class5-known-traceless}
\end{equation}
Then
\begin{align}
U_{5,0}
&=\mathcal G U_K\mathcal G^{-1}
 +(\partial_x\mathcal G)\mathcal G^{-1},
\label{eq:class5-known-U-map}\\
V_{5,0}
&=A_K^2\mathcal G V_K\mathcal G^{-1}
\label{eq:class5-known-V-map}
\end{align}
on the unit-spin constraint.  The spatial identity holds without the
constraint.  Before imposing the constraint, the difference in the time
identity is proportional to
$\widetilde{\boldsymbol S}^{\mathsf T}\widetilde{\boldsymbol S}-1$.

The scalar parts are
\begin{equation}
a_U:=\frac12\tr U_5,
\qquad
a_V:=\frac12\tr V_5,
\label{eq:class5-known-scalar-parts}
\end{equation}
and, on the unit-spin constraint,
\begin{align}
a_U&=\frac{1}{4\lambda}+\frac{\alpha_5}{4}S^+,
\nonumber\\
a_V&=\frac{\ii\alpha_5}{2}
(S^+\partial_xS^3-S^3\partial_xS^+)
+\frac{\ii\alpha_5^2}{4}(S^+)^2.
\label{eq:class5-known-scalars}
\end{align}
Using the residual $E^{(5)}_+$ defined in
Eq.~\eqref{eq:class5-residual-plus},
\begin{equation}
\partial_ta_U-\partial_xa_V
=\frac{\alpha_5}{4}E^{(5)}_+.
\label{eq:class5-known-scalar-curvature}
\end{equation}
The scalar sector is locally flat on the Class 5 equations of motion and can
be removed on shell by a local scalar auxiliary transformation.

For a periodic spin with circumference $\ell$,
\begin{equation}
\widetilde{\boldsymbol S}(x+\ell,t)
=\exp\!\left(\frac{\alpha_5\ell}{2}M_5\right)
\widetilde{\boldsymbol S}(x,t),
\qquad
G(x+\ell,\lambda)G(x,\lambda)^{-1}
=\Id_2-\frac{\alpha_5\ell}{2}E_{12}.
\label{eq:class5-known-boundary-twist}
\end{equation}
The field and gauge transformations are not generically periodic, so the
local correspondence does not identify periodic monodromies or global
boundary sectors.  Sections~\ref{sec:class5-time-lax} and
\ref{sec:class5-eom} give the direct finite-lattice derivation of the Class 5
time component.
